\section{Why Legal Queries Fail}\label{sec:preliminary}

\paragraph{Generated search text.}
RAG combines a generator with an evidence store \citep{lewis2020rag}. Its
legal promise is simple: retrieve the relevant rule or statute, then answer
from it. But the query and the corpus often use different language. HyDE
addresses this by generating a hypothetical document and embedding that text
for retrieval \citep{gao2023hyde}. Query2doc and related expansion methods
similarly use a language model to turn sparse query text into richer retrieval
handles \citep{wang2023query2doc}. For legal tasks, the useful generated text
is usually not a longer version of the question. It is a passage in the style
of the authority the retriever should find.

\paragraph{Legal RAG evaluation.}
Legal NLP benchmarks and corpora have standardized many legal-understanding
tasks \citep{guha2023legalbench,chalkidis2022lexglue,henderson2022pileoflaw},
and recent work has emphasized the risk of unsupported legal assistant answers
\citep{magesh2024hallucinations}. Legal RAG benchmarks such as LegalBench-RAG
and Legal RAG Bench argue for evaluating retrieval and answering together
\citep{pipitone2024legalbenchrag,butler2026legalragbench}. We follow that
discipline but keep the methods fixed: every row uses a predetermined retrieval
interface, and the analysis asks whether failure comes from missing evidence or
from using retrieved evidence poorly.

\paragraph{Two failure modes.}
BarExamQA and HousingQA expose different legal retrieval problems. BarExamQA
questions are fact patterns whose support passages often state compact legal
rules. HousingQA questions are jurisdiction-specific statutory questions over
a national corpus. The first setting tests whether generated legal passages
can translate facts into doctrine. The second tests whether the corpus has
been scoped to the right jurisdiction before method comparisons begin.
